\clearpage
\section{\MakeUppercase{Apresentação do Projeto Pedagógico do Curso}}
~ % Utilizado para saltar linha com recuo

O presente documento tem por finalidade apresentar e fundamentar a proposta pedagógica do \NOMECURSO do Instituto Federal de Educação, Ciência e Tecnologia do Ceará (IFCE), \textit{Campus} Maracanaú, o qual é ofertado na modalidade presencial. A estrutura deste Projeto Pedagógico de Curso (PPC) está amparada na Lei de Diretrizes e Bases da Educação Nacional (LDB), Lei nº 9.394/1996, e nas normativas legais vigentes em âmbito nacional e institucional que orientam a oferta e a organização dos cursos de graduação, com ênfase nas diretrizes específicas aplicáveis aos cursos de bacharelado.

A concepção inicial do \textbf{\NOMECURSO} foi elaborada pela comissão instituída pela Portaria \textbf{\PORTARIACOMISSAO} e homologada pela Resolução \textbf{\RESOLUCAOAPROVACAO}. Posteriormente, o PPC foi atualizado pela Resolução \textbf{\RESOLUCAOAPROVACAONOVA}, resultado do trabalho conjunto e articulado entre o Núcleo Docente Estruturante (NDE) e o Colegiado de Curso. 

No decorrer de seu desenvolvimento institucional, o PPC passou por revisões sucessivas que ajustaram pré-requisitos, cargas horárias, conteúdos e a organização curricular, além de promoverem a consolidação de laboratórios e a atualização de documentos acadêmicos. Essas revisões expressam um processo contínuo de aperfeiçoamento, orientado pela necessidade de manter o curso alinhado às orientações institucionais e às diretrizes nacionais para a formação em Engenharia, assegurando a coerência da proposta formativa e a atualização permanente da estrutura curricular..

A última revisão do PPC, apresentada nesta versão, foi conduzida pelo NDE, instituído pela Portaria \PORTARIANDEATUAL. Seu propósito central foi garantir a conformidade do curso com as diretrizes institucionais e nacionais vigentes, assegurando a coerência entre a proposta formativa e as normativas que regem os cursos de Engenharia.

A fundamentação legal completa para este PPC encontra-se detalhada na Seção \ref{FUNDAMENTAÇÃO LEGAL}. Contudo, para referência imediata, destaca-se que esta revisão foi orientada, sobretudo, pela Resolução CNE/CES nº 2, de 24 de abril de 2019 (\citeonline{cne19}), que estabelece as Diretrizes Curriculares Nacionais para os cursos de Engenharia, pela Resolução CNE/CES nº 1, de 26 de março de 2021 (\citeonline{cne21}), que atualiza e complementa essas diretrizes, e pela Resolução CONSUP/IFCE nº 259, de 8 de janeiro de 2025 (\citeonline{res259}), que institui o Manual de Normalização de Projetos Pedagógicos dos Cursos de Engenharia do IFCE. Também foram considerados o Guia de Curricularização das Atividades de Extensão nos Cursos Técnicos, de Graduação e Pós-Graduação do IFCE (\citeonline{GuiaCurricExt2023}) e o Instrumento de Avaliação de Cursos de Graduação do INEP (\citeonline{InstrAval2017}), como referenciais para alinhamento institucional e conformidade avaliativa.


A incorporação desses referenciais assegura a conformidade do PPC com as exigências legais, acadêmicas e formativas relativas à formação em Engenharia, conforme detalhado nas seções subsequentes deste documento.
